% from main.tex
%%%%%%%%%%  scan idx 140  |  folio 127  %%%%%%%%%%
pour \(i<b\), \(F'^b=Z^b(F)\), \(F'^i=0\) pour \(i>b\). Mais \(Z^b(F)\) admet une résolution finie droite par des objets de \(\mathcal C_0\) donc est localement objet de \(\mathcal C_0\), d'où l'assertion.

\textbf{Corollaire 4.14.} Soient \(F\in\Ob\mathcal C_S\) et \(n\in\mathbb N\). Conditions équivalentes :
\begin{enumerate}[label=(\roman*)]
\item \(\perfamp(F)\subset[-n,0]\);
\item \(F\) admet localement une résolution gauche de longueur \(n\) par des objets de \(\mathcal C_0\), i.e. l'ensemble des \(X\) au-dessus de \(S\) tels qu'il existe une suite exacte
\[
  0\longrightarrow L^{-n}\longrightarrow L^{-n+1}\longrightarrow \cdots
  \longrightarrow L^0\longrightarrow F_X\longrightarrow 0,
\]
avec \(L^i\in\Ob\mathcal C_{0,X}\) pour tout \(i\), est un raffinement de \(S\).
\end{enumerate}

\emph{Preuve.} C'est une conséquence immédiate de (4.7) et (4.13).

\textbf{Définition 4.14.1.} \emph{Nous dirons que \(F\) est de dimension parfaite \(\leq n\), et nous écrirons \(\perfdim(F)\leq n\), si \(F\) vérifie les conditions équivalentes de (4.14). Nous appellerons dimension parfaite de \(F\) le plus petit entier \(n\) tel que \(\perfdim(F)\leq n\) (s'il n'existe pas un tel entier, nous dirons que \(F\) est de dimension parfaite infinie).}

Quand \(S\) est le site ponctuel et que \(\mathcal C_0\) est la sous-catégorie pleine de \(\mathcal C\) formée des objets projectifs, la dimension parfaite n'est autre que la dimension projective usuelle.

\textbf{Proposition 4.15.} \emph{Soient \(a,b,n\in\mathbb Z\) (\(a\leq b\), \(n\geq0\)), et soit \(F\in\Ob\Dpx(\mathcal C)\) tel que \(F^i=0\) (resp. \(\Ho^i(F)=0\)) pour \(i\notin[a,b]\). Si, pour tout \(i\in[a,b]\), \(F^i\) (resp. \(\Ho^i(F)\)) est de dimension parfaite \(\leq n+(i-a)\), alors \(F\) est d'amplitude parfaite contenue dans \([a-n,b]\). Par suite, si pour tout \(i\in[a,b]\), \(F^i\) (resp. \(\Ho^i(F)\)) est parfait, \(F\) l'est aussi.}

